\section{The Anti-Christ \& the Liberal-Democratic Order}

In an interesting turn of phrase, St. Paul indicates in Thessalonians 2 that the Anti-Christ will come, not as an idol, but as an iconoclast:

\begin{quotex}
He who opposes and exalts himself against everything that is called God and religion, just as he will sit in the Temple of God, as God, and will show concerning himself as if he is God...\footnote{\url{http://bible.cc/2\_thessalonians/2-4.htm}}
\end{quotex}

Now this is an interesting idea, because in our modern world (particularly the Christian embers which sustain the liberal religion), there is no greater sin than idolatry or ``false Gods''. Ever since Calvin wrote the Institutes\footnote{\url{http://www.pravoslavieto.com/inoverie/protestantism/reformation\_iconoclasm.htm}}, Western mankind has had a tremendously hyperactive spiritual defense system against anyone who attempts to keep elements of ``paganism'' or ``idol-worship'' within their
\emph{weltanschaung}. Within the secular liberal democratic order which dominates public society, there is no greater crime than some hierarchic, ``obsolete'' ``authority'' (double quotes intentional) offending or even worse transgressing in some way against the hapless and defenseless individual. The appearance of this theme first makes its way in the West in the play Antigone\footnote{\url{https://en.wikipedia.org/wiki/Antigone}}, in which there is a tragic conflict between law of the \emph{polis} and the piety of a daughter who is burying her family. This tension is later exacerbated and intensified after the relapse of the Middle Ages and the expiration of the Baroque into the ``classical liberal\footnote{\url{https://en.wikipedia.org/wiki/Classical\_liberalism}}'' movement, which is a highly spiritual exaltation and subtle secularization of medieval Christian mores. Although a strong case can be made that classical liberalism is ``not for export'' to a thoroughly secular age (and requires a traditional society within which to subsist), I will not assume this for purposes of beginning a discussion; we can instead ``credit'' secularism with classical liberalism's goods, such as they are.

Let us suppose that classical liberalism is to be credited with much good -- is it separable from anti-traditional views of authority? Or put another way, what positive good does Lord Acton have to replace the Middle Ages? Worse, what positive good will Fukuyama\footnote{\url{https://en.wikipedia.org/wiki/The\_End\_of\_History\_and\_the\_Last\_Man}} propose to replace classical liberalism? How will the Great Society protect us from an Anti-Christ?

These are not idle debating points posed for rhetorical exercise -- many classical liberals today are concerned about the rising tide of Socialism and the power of demagogues to bring it to pass -- but they are in a difficulty, since it is not clear what separates one iconoclasm from the next.

Because the Anti-Christ (insists St. Paul) will ot ncome directly, but will proceed by way of iconoclasm. Moderns will attack ``privilege'' (economic idolatry) and the post-moderns will assault ``property'' (economic advantage resulting from remnants and dregs of aforementioned idolatry), all in the name of a God, Who is to come (one might add). Anti-Christ will subtly substitute itself into the place of God through this very process, inserting its own vision of ``the Good'' and ``God'' in place of the residual values of the old which it uses as a cloak and vehicle to infiltrate that antiquated order. Here is a website (for instance\footnote{\url{https://onecosmos.blogspot.com/}}) which generally argues using classical liberal assumptions, and is twisting in tenterhooks, right now.

My question would be this: what separates the very philosophy of classical liberalism\footnote{\url{http://distributistreview.com/mag/2012/08/lord-acton-tends-to-corrupt/}} (pick its best advocates, such as John Adams or Lord Acton or Vinet\footnote{\url{https://en.wikipedia.org/wiki/Alexandre\_Vinet}}) from the same mechanism or tendency? Why draw the line where they want it to lie? What reason can they produce, for this insistence? Because Western liberalism has consistently operated on the basis of opposing ``idolatry'' and substituting a secular vision of life in the place of ``older ways'' and false gods. This has been true since the Reformation, at the least. So why is classical liberalism not itself a servant unwilling of Anti-Christ?

For in St. Paul's prophecy, the Anti-Christ \emph{will not openly proclaim himself as God}, but will rather prepare the path for this by seemingly destroying little idols and objects of worship, a process which would necessarily involve the death of Tradition and its reverence for antique mores and thoughts, since liberalism is incapable of distinguishing true or false authority, particularly on an antique or vestigial basis. In fact, Anti-Christ (we may presume) may not ever openly proclaim itself as God, but will allow the process of substitution \emph{to itself be the end} -- The Means will be the End. What better description of modern Times\footnote{\url{http://fabiusmaximus.com/2012/09/26/russia-pussy-riot-43530/}} could one find, in which economic activity and political activity and technological progress will themselves have no \emph{raison d'etre} except the devouring of idols and the perpetuation of the endless circle of the Serpent, without any vertical lift to the spiral at all? At the round table of the Serpent, capitalist, liberal, and revolutionary sit together.

In this account, classical liberalism bites its own tail of Socialism. As a matter of fact, Hillaire Belloc has already given a clinically diagnostic account of this very process in \emph{The Servile State}\footnote{\url{http://archive.org/details/servilestate00belluoft}}. It is not however necessary to trust Belloc's argument (which is very tightly reasoned) -- one could simply note the prophecy of St. Paul. This would be enough to warn everyone that the modern iconoclasts are quite likely as spiritually empty as their dark master, Satan himself, and are constituting Anti-Christ.

Luther had his chance\footnote{\url{http://www.pravoslavieto.com/inoverie/protestantism/luther/luther\_had\_his\_chance.htm}}. The West has had its chance. It has instead opted for a virulent intensification of the fever of liberalism in the form of what we have today; however, unlike in Sophocles' times, there is no ``tragedy'' involved in the disobedience of the individual to the laws of the state. No one is in serious danger of being ``walled up'' or burnt at the stake (unless one is a ``reactionary''). Rather, the liberal democracies are exploiting the very chaos which they have engendered, in return for more power, money, and the opportunity to slake their passions. The war against Idols became a war against Nature, which became a war against Man himself, as was pointed out long ago:

\begin{quotex}
Almost as a postscript to the heavenly warning issued at Fatima in 1917, Saint Maximilian Kolbe, two years later, reviewed the three Great Evils of the latter times, noting: ``In 1517, the Protestants rebelled against the Church; in 1717, the Freemasons rebelled against Christ; and, in 1917, the Communists rebelled against God.'' In a single sentence the Polish martyr had exposed common origin and natural succession of each of these Apocalyptic nightmares...
\end{quotex}


C.M. Hofbauer\footnote{\url{http://catholicism.org/clement-maria-hofbauer.html}}


\flrightit{Posted on 2012-09-28 by Logres}

\begin{center}* * *\end{center}

\begin{footnotesize}
\begin{sffamily}
\texttt{Logres on 2012-09-29 at 21:07 said:}

I should note that the article should be clearer that, under liberal-democracy, in fact, there is in fact no guarantee possible against an Anti-Christ, because there is no Authority. And they would admit this (Lord Acton). In Biblical language, there is no ``one who restrains'' and ``stands in the way''. Gornahoor and Cologero have made clear before that the Ideal has often been outright perverted into Anti-Christ, but that is not the pattern we are dealing with in our age -- Satan has dispensed with the Ideal (which he discredited in principle and absolutely) -- he now has to operate only with false ``ideals'' which can be replaced with successively weaker imitations at his whim when purpose is exhausted. The new replacements have spiritual DNA which degenerate obviously and openly.

\hfill

\texttt{lckr1991 on 2012-09-30 at 00:50 said:}

Great article.

\hfill

\texttt{Shining Emperor on 2013-01-22 at 07:56 said:}

The Constitution UNmade law, Americans. Only the far-seeing know this, but bloodbaths of hecatombs of unoffending herded children slaughtered every month by grown men, constitutionally psychopathic (i.e., American), these ``signs'' of the times shall increasingly befall us who yet desire Order and side with the Catechon against the Anarch.


\hfill

\end{sffamily}
\end{footnotesize}

